\begin{textsample}{POS Dim 3 – ma – Score 8.00 – ma/ma\_inf\_8.txt}  \label{ex:f3_pos_170}
A avaliação na Educação \textbf{Infantil} A avaliação é uma atividade importante em qualquer processo educativo.
Seu objetivo deve ser conhecer melhor em que estágio as \textbf{crianças} estão e identificar de que modo é possível atuar de maneira mais efetiva para assegurar o seu pleno desenvolvimento.
Assim, constitui um recurso pedagógico adicional para os professores e coordenadores pedagógicos ou diretores de unidades de Educação \textbf{Infantil}.
A avaliação na Educação \textbf{Infantil} deve incidir diretamente no planejamento das atividades \textbf{diárias} promovidas pelo professor junto às \textbf{crianças}, devendo subsidiar elementos que ampliem as aprendizagens e experiências apresentadas por elas, contribuindo também para suas manifestações, \textbf{desejos} e necessidades.
Vale ressaltar que a avaliação nesta etapa da Educação Básica possui algumas especificidades que devem ser observadas.
A mais importante delas é o fato de que os procedimentos avaliativos não devem ter caráter de classificação ou promoção das \textbf{crianças} de uma fase a outra.
Sua finalidade maior é educativa.
A LDB de 1996 estabelece que a Educação \textbf{Infantil} precisa ser organizada com base em algumas regras comuns, entre elas “ avaliação mediante \textbf{acompanhamento} e \textbf{registro} do desenvolvimento das \textbf{crianças}, sem o objetivo de promoção, mesmo para o acesso ao Ensino Fundamental ” ( incluído pela Lei no 12.796/13 ).
Assim, o foco da avaliação está voltado para o \textbf{registro} de todas as conquistas, avanços, dificuldades e desafios enfrentados pelas \textbf{crianças}, com a finalidade de observar seu progresso no processo de ensino e aprendizagem.
Importante, ainda, considerar o disposto nas Diretrizes Curriculares Nacionais para a Educação \textbf{Infantil} de 2009, que reitera que “ as instituições de Educação \textbf{Infantil} devem criar procedimentos para \textbf{acompanhamento} do trabalho pedagógico e para avaliação do desenvolvimento das \textbf{crianças}, sem objetivo de seleção, promoção ou classificação ”.
Conforme detalhado no documento, essa avaliação deve garantir : I – a observação crítica e criativa das atividades, das \textbf{brincadeiras} e interações das \textbf{crianças} no cotidiano ; II – utilização de múltiplos \textbf{registros} realizados por \textbf{adultos} e \textbf{crianças} ( relatórios, fotografias, desenhos, álbuns etc. ) ; III – a continuidade dos processos de aprendizagens por meio da criação de estratégias adequadas aos diferentes momentos de transição vividos pela \textbf{criança} ( transição casa/instituição de Educação \textbf{Infantil}, transições no interior da instituição, transição creche/pré-escola e transição pré-escola/Ensino Fundamental ) ; IV – documentação específica que permita às famílias conhecer o trabalho da instituição junto às \textbf{crianças} e os processos de desenvolvimento e aprendizagem da \textbf{criança} na Educação \textbf{Infantil} ; V – a não retenção das \textbf{crianças} na Educação \textbf{Infantil}.
Essa passagem mostra a importância de, por um lado, propiciar experiências educativas cheias de significados para as \textbf{crianças} e, por outro, registrar seu desenvolvimento ao longo das atividades, apontando desafios e progressos e, sobretudo, estimulando suas potencialidades.

% matched lemmas: acompanhamento, adulto, brincadeira, criança, desejo, diário, infantil, registro
\end{textsample}
